% ===================================================================
%  BAGAN No. 15 — Pasar. --- Bahan makanan.
%
%  Traduction, faite en 2026, en indonesien d'aujourd'hui, sur l'ido
%  de texto/io. Regles de la colonne : voir 10-bagan-01.tex. Une seule
%  numerotation.
%
%  pasar EST LE MEME MOT QUE بازار, ET C'EST LA TROISIEME COLONNE OU
%  IL SE PRESENTE. Le persan bāzār est passe en ourdou par la terre,
%  en indonesien par la mer avec les marchands venus de l'ouest, et en
%  ido par l'italien et le francais. Trois routes, un seul mot, et il
%  nomme dans les trois colonnes le sujet du tableau. Le releve avait
%  note au tableau 16 ourdou que « La Bazaro » de l'ido est un mot
%  persan rendu a sa langue ; ici il est rendu une seconde fois, par
%  l'autre bout du monde.
%
%  LE MARCHE EST L'ENDROIT OU LE PARTAGE DES COUCHES SE VOIT LE MIEUX,
%  ET IL SUIT LA BOTANIQUE COMME AU TABLEAU 8.
%
%   * Le potager est NEERLANDAIS d'un bout a l'autre : wortel la
%     carotte, buncis les haricots, kol le chou, kembang kol le
%     chou-fleur, prei le poireau, seledri le celeri, selada la salade
%     (celle-ci portugaise). Ce sont les legumes que la Compagnie a
%     introduits dans les jardins de Java, et ils ont garde le nom de
%     ceux qui les ont plantes.
%   * Les poissons sont europeens et s'empruntent presque tous :
%     haring, sarden, kod, trout. La mer d'Indonesie est l'une des
%     plus poissonneuses du monde, mais ce ne sont pas ces
%     poissons-la.
%   * Et le fond est malais : beras le riz, pisang la banane, nanas
%     l'ananas, telur les œufs, jamur les champignons, bawang putih
%     l'ail, bawang merah l'oignon, bayam les epinards, daging la
%     viande, ikan le poisson, kepiting le crabe, udang la crevette.
%
%  MAIS C'EST LA BOUTIQUE DE L'EPICIER (32) QUI PORTE LE TABLEAU.
%  toko rempah, « la boutique aux epices » : le poivre, le clou de
%  girofle et la noix de muscade pour lesquels l'Europe a arme des
%  flottes pendant trois siecles viennent de cet archipel et de nulle
%  part ailleurs. merica, cengkih, pala sont des mots malais que
%  l'Europe est allee chercher a la source. Sur une planche belge de
%  1926 ils sont ranges parmi les articles ordinaires d'une epicerie
%  de quartier, a cote du savon et du sucre en pain. La colonne les
%  ecrit sans un mot de plus : c'est au lecteur de voir que le tableau
%  est chez lui.
%
%  QUATRIEME \VUgras{} VIDE DE LA COLONNE (apres les tableaux 5, 9 et
%  10), et toujours a la meme place : un renvoi qui suit un terme dont
%  le groupe gras est deja ferme. La main ouvre un second groupe pour
%  « porter » l'exposant, alors que l'exposant se porte tout seul.
%  Quatre fois sur quinze tableaux, c'est la faute propre de cette
%  colonne, et kolonoj.py la releve chaque fois : le verificateur n'a
%  pas eu a etre invente pour elle, il existait deja.
% ===================================================================

%%K t15-apar-1 apar
\VUsaut{4.0mm}
\VUtitre{15.0pt}{60}{BAGAN No. 15}
\VUsaut{3.0mm}

%%K t15-tit-1 sub
\VUpk{11.4pt}{40}{Pasar. --- Bahan makanan.}
\VUsaut{3.0mm}

%%K t15-01-1 p
1. --- Sebagian besar pedagang yang menyediakan bagi kita barang-barang
yang diperlukan untuk makanan kita berkumpul di bawah \VUgras{aula}
\textsuperscript{(1)} \VUgras{pasar tertutup} atau di jalan-jalan
sekitarnya.

%%K t15-02-1 p
2. --- \VUgras{Penjual daging babi} \textsuperscript{(2)}, yang menjual
\VUgras{ham} \textsuperscript{(3)}, \VUgras{sosis darah}
\textsuperscript{(4)}, \VUgras{sosis kecil} \textsuperscript{(5)},
\VUgras{sosis babat} \textsuperscript{(6)} dan \VUgras{lemak babi}
\textsuperscript{(7)}, berdiri di dekat \VUgras{penjual mentega dan
keju} \textsuperscript{(8)}, yang pajangannya penuh dengan \VUgras{keju
Belanda} \textsuperscript{(9)} berkulit merah, \VUgras{keju Camembert}
\textsuperscript{(10)}, \VUgras{keping-keping keju Gruyère}
\textsuperscript{(11)} dan \VUgras{bongkah mentega}
\textsuperscript{(12)} yang segar.

%%K t15-03-1 p
3. --- Di depannya, seorang \VUgras{penjual sayur} \textsuperscript{(13)}
menawarkan kepada para pembelinya \VUgras{tomat} \textsuperscript{(14)}
yang bagus, \VUgras{kembang kol} \textsuperscript{(15)},
\VUgras{selada} \textsuperscript{(16)}, \VUgras{kol}
\textsuperscript{(17)}, \VUgras{labu} \textsuperscript{(19)},
\VUgras{melon} \textsuperscript{(18)} dan bahkan \VUgras{jamur}
\textsuperscript{(20)}. Tetapi seorang \VUgras{pesuruh bir,} yang
menghentikan \VUgras{kereta antarnya} \textsuperscript{(21)} yang
sarat dengan \VUgras{tong-tong bir} \textsuperscript{(22)} tepat di
depan pajangannya, menghalangi para pembelinya mendekat. Seorang
tukang kebun perempuan membawakan kepadanya sekeranjang besar
\VUgras{artisok} \textsuperscript{(23)}.

%%K t15-tit-2 sub
\VUsaut{4.0mm}
\VUpk{10.2pt}{40}{Menjual dan membeli.}
\VUsaut{3.0mm}

%%K t15-04-1 p
4. --- Di pasar itu suasananya ramai, sebab waktu \VUgras{lelang}
\textsuperscript{(24)} telah tiba. \VUgras{Para ibu rumah tangga}
\textsuperscript{(26)}, yang datang ke sana untuk berbelanja, berhenti
ketika lewat di depan kios \VUgras{penjual babat} \textsuperscript{(25)}
untuk membeli \VUgras{babat} atau \VUgras{kepala anak sapi.}

%%K t15-05-1 p
5. --- Seorang \VUgras{juru masak} \textsuperscript{(27)} yang gemuk
menawar seekor \VUgras{tuna} yang sangat bagus pada \VUgras{penjual
ikan} \textsuperscript{(28)}, yang menaikkan harganya sesuka hati. Ia
menerima kiriman \VUgras{belut}, \VUgras{ikan lidah}, \VUgras{ikan
trout} dan \VUgras{ikan mas} dalam jumlah besar. \VUgras{Ikan kodnya}
dan \VUgras{kerangnya} sangat segar.

%%K t15-tit-3 sub
\VUsaut{4.0mm}
\VUpk{10.2pt}{40}{Barang murah dan barang mahal.}
\VUsaut{3.0mm}

%%K t15-06-1 p
6. --- \VUgras{Penjual unggas} \textsuperscript{(29)}, yang
berjualan di dekatnya, sangat jujur. Bila ia menjual \VUgras{kalkun,}
\VUgras{angsa,} \VUgras{itik} atau sekadar \VUgras{anak ayam,} ia hanya
meminta harga yang pantas.

%%K t15-07-1 p
7. --- Di sudut alun-alun, di lantai satu, sejak beberapa waktu
menetap seorang \VUgras{penjual kain} \textsuperscript{(30)}, yang di
tempatnya para pembeli perempuan selalu pasti menemukan pilihan
\VUgras{taplak meja,} \VUgras{serbet,} \VUgras{celemek} dan \VUgras{lap}
yang sangat bagus. Di lantai yang sama, seorang \VUgras{tukang pisau}
\textsuperscript{(31)} membuka bengkelnya. Ia mengasah \VUgras{pisau}
para penjual perempuan di aula itu dan membuat untuk para tukang kebun
\VUgras{sabit} dari \VUgras{baja} yang paling keras.

%%K t15-tit-4 sub
\VUsaut{4.0mm}
\VUpk{10.2pt}{40}{Toko rempah.}
\VUsaut{3.0mm}

%%K t15-08-1 p
8. --- Di bawah bengkelnya, sejak beberapa waktu dibuka sebuah toko
bahan makanan yang besar. Ia bukan lagi \VUgras{toko rempah}
\textsuperscript{(32)} yang sederhana dan tidak penting seperti dahulu,
yang hanya menjual barang sehari-hari : \VUgras{teh}
\textsuperscript{(33)}, \VUgras{cokelat} \textsuperscript{(34)},
\VUgras{gula batu} \textsuperscript{(37)} dan \VUgras{sabun}
\textsuperscript{(38)}. Di sana orang menjual segala barang :
\VUgras{biskuit renyah,} \VUgras{makanan kaleng}
\textsuperscript{(35)}, \VUgras{minuman keras manis}
\textsuperscript{(36)}, \VUgras{rum,} \VUgras{konyak,}
\VUgras{kirsch,} \VUgras{minuman adas} dan \VUgras{curaçao.} Di sana
orang menemukan daging buruan : \VUgras{kelinci hutan}
\textsuperscript{(39)}, \VUgras{burung anis} \textsuperscript{(40)},
\VUgras{ayam hutan} \textsuperscript{(41)}, \VUgras{burung puyuh}
\textsuperscript{(42)}, \VUgras{itik liar} \textsuperscript{(43)} atau
\VUgras{burung pegar} \textsuperscript{(44)}, semudah buah dari negeri
jauh seperti \VUgras{pisang} \textsuperscript{(46)} dan \VUgras{nanas.}

%%K t15-09-1 p
9. --- \VUgras{Pelayan} \textsuperscript{(45)} toko rempah itu sangat
suka menolong dan selalu siap memberikan apa yang diminta :
\VUgras{selai} \textsuperscript{(47)} kismis atau kuinsi,
\VUgras{lemak} \textsuperscript{(48)}, \VUgras{acar ketimun}
\textsuperscript{(49)} atau \VUgras{sarden} \textsuperscript{(50)}
asin.

%%K t15-09-2 p
Hal itu sangat memudahkan \VUgras{para pembeli perempuan}
\textsuperscript{(51)}, yang menemukan, di samping barang sehari-hari,
\VUgras{buah-buah pertama} yang paling langka dan buah yang paling
lezat : \VUgras{buah pir} \textsuperscript{(52)} yang berair,
\VUgras{prem} \textsuperscript{(53)}, \VUgras{anggur}
\textsuperscript{(54)}, \VUgras{buah persik} \textsuperscript{(55)},
\VUgras{badam} \textsuperscript{(56)} dan \VUgras{buah walnut}
\textsuperscript{(57)}.

%%K t15-tit-5 sub
\VUsaut{4.0mm}
\VUpk{10.2pt}{40}{Para pembeli.}
\VUsaut{3.0mm}

%%K t15-10-1 p
10. --- \VUgras{Beras} \textsuperscript{(58)} dan \VUgras{miju-miju}
\textsuperscript{(59)} ada di dekat peti \VUgras{ikan haring}
\textsuperscript{(60)} asap ; \VUgras{kentang} \textsuperscript{(61)}
dan \VUgras{buncis} \textsuperscript{(63)} bertetangga dengan
\VUgras{apel} \textsuperscript{(62)}, \VUgras{buah ara}
\textsuperscript{(64)}, \VUgras{prem kering} \textsuperscript{(65)} dan
\VUgras{buah hazel} \textsuperscript{(66)}. Di toko itu orang menemukan
\VUgras{telur} \textsuperscript{(67)} segar dan \VUgras{zaitun}
\textsuperscript{(68)} ; karena itu \VUgras{keranjang}
\textsuperscript{(70)} dan \VUgras{jaring belanja}
\textsuperscript{(73)} terisi dengan sangat cepat.

%%K t15-11-1 p
11. --- \VUgras{Kopi} \textsuperscript{(72)} toko yang bagus itu
memperoleh nama baik yang pantas di kalangan \VUgras{para pembantu}
\textsuperscript{(69)} dan \VUgras{para juru masak,} sebab
\VUgras{pedagang rempah} \textsuperscript{(71)} tua itu sendiri yang
menyangrainya dengan sangat cermat.

%%K t15-tit-6 sub
\VUsaut{4.0mm}
\VUpk{10.2pt}{40}{Tontonan.}
\VUsaut{3.0mm}

%%K t15-12-1 p
12. --- Tidak semua orang datang ke pasar untuk berbelanja. Dalam lalu
lalang yang menarik di gang yang menuju aula itu, orang melihat
manusia yang paling bermacam-macam. Di dekat seorang \VUgras{pesuruh
tukang daging} \textsuperscript{(75)} yang ramah, yang mendorong di
depannya sebuah \VUgras{gerobak dorong} \textsuperscript{(74)}, inilah
dua prajurit yang sedang cuti, yang sangat bangga memperlihatkan
seragam mereka yang cemerlang : seorang \VUgras{husar}
\textsuperscript{(76)} dengan \VUgras{dolman} birunya dan \VUgras{topi
berbulunya,} dan seorang \VUgras{kopral zuaf} dengan lengan bajunya
yang menggembung dan kepalanya yang bertutup \VUgras{tarbus.}

%%K t15-13-1 p
13. --- \VUgras{Tukang roti} \textsuperscript{(80)}, yang berdiri di
ambang \VUgras{toko rotinya} \textsuperscript{(78)}, baru saja
menguleni adonan \VUgras{rotinya} \textsuperscript{(79)} dan
mengeluarkan dari tungku \VUgras{roti-roti bulan sabit}
\textsuperscript{(81)}, \VUgras{roti gandum halus}
\textsuperscript{(82)} dan \VUgras{roti bundar} \textsuperscript{(83)}
yang ada di etalase.

%%K t15-14-1 p
14. --- Di atas toko roti itu tinggal seorang \VUgras{penjahit kain}
\textsuperscript{(84)} yang cakap, yang mempekerjakan beberapa
\VUgras{penyulam} \textsuperscript{(85)}. Di sebelahnya tinggal seorang
\VUgras{tukang cuci dan setrika} \textsuperscript{(86)}, yang mengkanji
\VUgras{kain} \textsuperscript{(88)} sesudah mencuci dan
mengeringkannya, lalu melicinkannya dengan \VUgras{setrika}
\textsuperscript{(87)}.

%%K t15-tit-7 sub
\VUsaut{4.0mm}
\VUpk{10.2pt}{40}{Daging, ikan dan sayuran.}
\VUsaut{3.0mm}

%%K t15-15-1 p
15. --- Di \VUgras{tempat jagal} \textsuperscript{(89)}, yang terletak
di lantai dasar, orang menjual segala macam \VUgras{daging} —
\VUgras{daging domba} \textsuperscript{(90)}, \VUgras{daging sapi}
\textsuperscript{(94)} atau \VUgras{daging anak sapi}
\textsuperscript{(92)} — dalam bentuk \VUgras{paha domba, bistik,
daging panggang} dan \VUgras{kotelet.} \VUgras{Pesuruh tukang daging}
\textsuperscript{(93)} memotong-motong semua daging itu dan meletakkan
terpisah \VUgras{tulang-tulang sumsum} \textsuperscript{(96)}. Di pintu
tempat jagal itu ada seorang \VUgras{penjual tiram}
\textsuperscript{(97)} yang menjual \VUgras{tiram} \textsuperscript{(98)}.
Ia membukanya bila orang menginginkannya.

%%K t15-16-1 p
16. --- Semua pedagang itu membayar izin usaha atau retribusi tempat ;
karena itu \VUgras{polisi} \textsuperscript{(100)} tanpa ampun mengejar
\VUgras{para penjaja sayur} \textsuperscript{(99)} yang malang, yang
menyaingi mereka dengan mengedarkan di gerobak kecil mereka
\VUgras{wortel, lobak, lobak putih, prei} atau \VUgras{ikatan
asparagus, kacang polong segar, bayam, daun asam, selada air} dan
\VUgras{peterseli.}

%%K t15-tit-8 sub
\VUsaut{4.0mm}
\VUpk{10.2pt}{40}{Awas polisi!}
\VUsaut{3.0mm}

%%K t15-17-1 p
17. --- Anak yang menjual \VUgras{bawang putih} \textsuperscript{(101)}
dan \VUgras{bawang merah} tahu benar bahwa ia pun tidak boleh menjual
dagangannya di dekat pasar. Ia lari, dan sambil berlari ia berisiko
menjatuhkan keranjang \VUgras{kepiting,} \VUgras{udang karang kecil}
dan \VUgras{udang} milik \VUgras{penjual} \textsuperscript{(102)}
\VUgras{udang barong} dan \VUgras{lobster.}

%%K t15-18-1 p
18. --- Semua orang bergerak dan bersuara keras ; tetapi hal itu tidak
menghalangi orang mendengar dengan jelas suara \VUgras{tukang kaca}
\textsuperscript{(103)} yang berkeliling, maupun seruan \VUgras{penjual
arang} \textsuperscript{(104)} yang baru saja meninggalkan gerobaknya
sejenak untuk mengantarkan sebuah \VUgras{karung arang}
\textsuperscript{(105)}.

\VUsaut{6.0mm}
\VUfilet{22mm}
